\section{詳細解説}

\subsection{トルクとエネルギーの周波数依存}

人間大のモデルで、最適ループを回す周波数を 0.1〜2 Hz で変え、関節トルクの最大値と 90° 旋回に必要な正の仕事を求めた（図~\ref{fig:D5_torque_energy.png}）。
\begin{itemize}
\item 関節トルクは $f^{2.00}$ に比例した（付録~\ref{app:eom} の時間スケーリングのとおり）。
\item 90° 旋回に必要な正の仕事（回生なし）も $f^{2.00}$ に比例した。1 周あたりの仕事は $f^3$（パワー）÷ $f$（周回数/時間）$= f^2$ に比例し、必要な周回数は $f$ によらないためである。
  0.1 Hz では 20 J、2 Hz では 7900 J になる。\textbf{ゆっくり回すほど、エネルギーは大幅に少なくて済む}。
\end{itemize}

\DFIG{D5_torque_energy.png}{左：最適ループを周波数 $f$ で回すときの関節トルクの最大値（点線は持続トルクの目安）。右：90° 旋回に必要な正の仕事。どちらも $f^2$ に比例する。}

\subsection{相似拡大の全結果}

\DTAB{sim04_scaling}

\begin{example}[18 m 級の 90° 旋回を見積もる]
人間大（$s=1$）で、0.5 Hz の股関節トルクは 28.8 N\,m、持続トルクの目安は 120 N\,m なので、$f_{\max}=0.5\sqrt{120/28.8}=1.02$ Hz。
18 m 級（$s=10$）では、同じ周波数で必要なトルクが $10^5$ 倍、使えるトルクが $10^3$ 倍なので、$f_{\max}$ は $1/\sqrt{100}=1/10$ になり 0.102 Hz。
1 周 6.70° なので 90° には 13.4 周、$13.4/0.102 = 132$ s かかる。
\end{example}

\subsection{質量行列の構成}

順動力学では、重心系の質量行列 $M(q)$ を次のように組み立てた。
各剛体 $i$ の重心速度と角速度は一般化速度 $\bm u=(\dot\theta,\dot q)$ の 1 次式 $\bm v_i = J_{v,i}\bm u$、$\omega_i = J_{\omega,i}\bm u$ と書ける。
系重心の速度 $\bm v_c = J_c\bm u$（$J_c = \sum m_iJ_{v,i}/M$）を引いた相対速度を使うと
\begin{equation}
M(q) = \sum_i\left[m_i(J_{v,i}-J_c)^\top(J_{v,i}-J_c) + I_iJ_{\omega,i}^\top J_{\omega,i}\right].
\end{equation}
この $M$ の第 1 行が $(A, B)$ に一致することを数値で確認した（差 $10^{-15}$）。
コリオリ・遠心力の項は $\partial M/\partial q_k$ を中心差分で求めて組み立てた。

\subsection{検証の詳細}

トルクだけを入力にした順動力学で、0.5 Hz の最適ループを 2 周（4 s）積分した。
\begin{itemize}
\item 入力トルクは、運動学の軌道から逆動力学で求めたもの（フィードフォワードのみ、フィードバックなし）。
\item 結果の関節軌道は、目標の軌道と $2.6\times10^{-8}$° 以内で一致した（開ループでも 2 周の間ずれない）。
\item 2 周後の胴体の回転は $+13.4086^\circ$ で、運動学モデル（Sim03 のループ 1 周 $\times2$）と一致した。
\item 角運動量は $2\times10^{-9}$ 以内で保存され、運動エネルギーの変化と関節の仕事の差は $9\times10^{-4}$ J（台形積分の誤差）だった。
\end{itemize}

\section*{考察課題}
\begin{exercise}
関節に弾性要素（ばね）を入れ、周回運動の固有振動数をループの周波数に合わせると、必要な正の仕事はどう変わるか。図~\ref{fig:sim04_human_cycle.png} のパワーの正負に注目して考えよ。
\end{exercise}
\begin{exercise}
アクチュエータのトルク密度が質量に比例する（トルク上限 $\propto s^3$）のではなく、断面積に比例する（$\propto s^2$）とすると、$f_{\max}$ と 90° 旋回時間は $s$ のどのべきで変わるか。
\end{exercise}
